\section{Discussion and limitations}
\label{sec:discussion}

\paragraph{Cost, and where the value concentrates.}
The model's principal cost is attestation: machines can harvest candidates, but hardness and classification are human judgments, and judgments demand time from the teams with the least of it. Three considerations temper this. First, the marginal cost is lowest exactly where the value is highest---tier-1 services and the platform resources beneath them---and the metrics remain meaningful at partial coverage because coverage itself is measured ($\rdc$) and estimation error has a known sign (Cor.~\ref{cor:optimism}). Second, much of the knowledge is already being produced and discarded: every post-incident review and every drill generates edges and evidence that today evaporate into prose. Third, whether the residual cost is acceptable is not a matter for argument but for RQ1, which measures it. If tier-1 modelling costs an order of magnitude more than our pre-registered expectation, the model as proposed fails its own test, and we would regard that as a result worth publishing.

\paragraph{Staleness, incentives, and gaming.}
A Recovery Graph that drifts becomes a CMDB---confidently wrong. The design response is that staleness is itself first-class: evidence expires by construction, drift against inventory and traces is a continuously computed finding, and the metrics reward freshness rather than volume. The gaming analysis of Sec.~\ref{sec:metrics} is deliberately explicit; we do not believe metric suites without named failure modes, and we expect the separation of policy-setting from attestation (Sec.~\ref{sec:governance}) to matter more in practice than any formula. What the model cannot manufacture is an organisation that wants to know its recovery posture; it can only make wanting-to-know cheap and not-knowing visible.

\paragraph{Model expressiveness.}
Four boundaries of the present formalisation are worth stating plainly. (i) $\Lev$ is a total order; degraded modes that are incomparable (read-only vs.\ region-local) need the lattice generalisation, which the mechanisms admit but the paper does not develop. (ii) The base model is conjunctive: alternative recovery routes---restore from snapshot \emph{or} rebuild from source---require and/or structure, at which point necessity analysis becomes dominator-based (Sec.~\ref{sec:blocking}) and plan choice becomes optimisation. (iii) Durations are deterministic estimates; a stochastic extension (distributions on $d_v(\ell)$, yielding $\rto$ quantiles by simulation over $\mathcal{A}$) is straightforward but its calibration is not, and we prefer honest point estimates paired with coverage over uncalibrated intervals. (iv) Assumption A1 excludes mid-recovery regression; scenario S5 probes this boundary empirically, and a game-semantic treatment (recovery against an adversarial environment) is an intriguing theoretical direction.

\paragraph{Multi-team and multi-cloud scale.}
Nothing in the semantics is single-cluster, but federation raises governance questions the paper only gestures at: who attests cross-organisational edges (the payment provider's sandbox), how much of a vendor's recovery posture can be imported as third-party evidence, and how DORA-style obligations on ICT third parties \citep{eu2022dora} interact with edges that cross company boundaries. We suspect the graph is \emph{more} valuable at these boundaries---they are where tacit knowledge is thinnest---but the claim is untested.

\paragraph{Security of the artifact itself.}
A complete map of what must be recovered, in what order, with what credentials, is reconnaissance gold. The graph store therefore warrants the same protection level as the secret store it references, provenance on evidence limits fabrication, and the C6 out-of-band copy must be protected \emph{and} reachable---a genuine tension (availability of the plan vs.\ confidentiality of the plan) that deployments must resolve consciously, e.g.\ by tiering what the offline copy contains.

\paragraph{When not to use this.}
A five-service monolith with one database does not need a Recovery Graph; it needs a tested restore script and a printed runbook---which is to say, its recovery graph fits in a head. The model earns its overhead where no single head holds the closure: many teams, many stateful systems, platform layers with their own recovery structure, and auditors who ask questions that today take weeks of archaeology. Below that threshold, the checklist is the better technology.
